\begin{question}[index=11]
设
\[
\lim_{x\to+\infty}\bigl(\sqrt[3]{1+ax^3}+x\bigr)=0\text{，}
\]
则曲线 $y=\sqrt[3]{1+ax^3}$ 在 $(0,+\infty)$ 内的斜渐近线方程为\fillin。
\end{question}
\QuestionBreak

\begin{question}[index=12]
圆 $(x-1)^2+y^2=1$ 绕 $y$ 轴旋转一周所得旋转体的表面积为\fillin.
\end{question}
\QuestionBreak

\begin{question}[index=13]
设 $a^2+b^2=1$（$a<0,b>0$），则曲线 $y=x^2+ax$ 与直线 $y=bx$ 所围区域的面积的最大值为\fillin.
\end{question}
\QuestionBreak

\begin{question}[index=14]
设上半球面 $\Sigma:z=\sqrt{1-x^2-y^2}$，其密度 $\rho=z$，则曲面 $\Sigma$ 的质心坐标 $\overline z=$\fillin.
\end{question}
\QuestionBreak

\begin{question}[index=15]
设矩阵 $\vect A=\mat{1&0&3\\1&4&5\\0&0&2}$，$\vect B=\mat{1&2&1\\3&5&a\\2&5&7}$，矩阵 $\vect C$ 满足 $\vect A\vect C\vect A^*=\vect B$，$r(\vect C)=2$，其中 $\vect A^*$ 是 $\vect A$ 的伴随矩阵，则 $a=$\fillin.
\end{question}
\QuestionBreak

\begin{question}[index=16]
设随机变量 $X$ 与 $Y$ 相互独立，且均服从 $N(0,\sigma^2)$（$\sigma>0$），则 $P\{0\le\sqrt{X^2+Y^2}\le\sqrt2\sigma\}=$\fillin.
\end{question}
